% ---
% id: fig4b-gap-provenance
% title: 논문 Fig. 4(b) gap 값의 출처 확인
% status: blocked
% created: 2026-10-03
% next: 논문 그림의 J_PD = 0.010 meV 값을 아래 지문 값과 대조한다
% blocked-reason: 저장소 이력이 2026-09-29부터라 Fig. 4(b)가 legacy/scripts/4_Pseudo_Gap.py(×S)로 그려졌는지 판별할 수 없음
% resume-condition: 사용자가 그림 값과 지문 표를 대조해 알려 줄 때
% ---
\section{논문 Fig. 4(b) gap 값의 출처 확인}

% 근거: docs/nbcp/appendices/a-pseudo-goldstone-gap.tex
% 지문 값(meV, J = 0.010 meV): 올바른 식 Y/PD 0.0069, V/PD 0.0156, Y/Γ 7.3e-6, V/Γ 0.0020;
%   legacy 그대로 0.0024, 0.0047, 2.6e-7, 0.00020; legacy를 1/S로 고친 것 0.0098, 0.0189, 1.0e-6, 0.0008.
%   Γ 축(Y/Γ 28배 차이)이 가장 잘 구별된다.
% 출처: GOVERNMENT/Working-Pad/project-review/2026-10-01-park2026-gap-analysis.md §Fig. 4(b) provenance check
% 결론은 이 블록이 아니라 docs/nbcp/ 해당 장에 옮겨 적는다 (NBCP 원본은 main.tex)
